\documentclass[11pt,letterpaper]{article}

\usepackage[T1]{fontenc}
\usepackage[utf8]{inputenc}
\usepackage[spanish,es-nodecimaldot]{babel}
\usepackage{lmodern}
\usepackage[letterpaper,margin=2.35cm,headheight=22pt]{geometry}
\usepackage{xcolor}
\usepackage{graphicx}
\usepackage{booktabs}
\usepackage{tabularx}
\usepackage{array}
\usepackage{float}
\usepackage{tcolorbox}
\usepackage{fancyhdr}
\usepackage{hyperref}

\definecolor{ink}{HTML}{14282B}
\definecolor{petrol}{HTML}{0C5968}
\definecolor{copper}{HTML}{C7643B}
\definecolor{mist}{HTML}{F3F7F7}
\definecolor{line}{HTML}{CCD9D9}
\definecolor{muted}{HTML}{587074}

\hypersetup{
  colorlinks=true,
  linkcolor=petrol,
  citecolor=petrol,
  urlcolor=copper,
  pdftitle={Text-to-SQL con QLoRA},
  pdfauthor={Wilson Alejandro Calderón Argueta y Pablo Daniel Barillas Moreno}
}

\pagestyle{fancy}
\fancyhf{}
\lhead{\small\color{muted}Deep Learning y Sistemas Inteligentes}
\rhead{\small\color{muted}Proyecto 3}
\cfoot{\small\color{muted}\thepage}
\renewcommand{\headrulewidth}{0.4pt}
\renewcommand{\headrule}{\hbox to\headwidth{\color{line}\leaders\hrule height \headrulewidth\hfill}}

\newtcolorbox{finding}[1]{
  colback=mist,
  colframe=petrol,
  boxrule=0.6pt,
  arc=2mm,
  title=\textbf{#1},
  coltitle=ink
}

\newcommand{\pending}{\textcolor{copper}{\textbf{Pendiente de medición}}}

\begin{document}
\pagenumbering{gobble}

\begin{titlepage}
  \color{ink}
  \vspace*{1.3cm}
  {\small UNIVERSIDAD DEL VALLE DE GUATEMALA\par}
  \vspace{0.35cm}
  {\color{petrol}\noindent\rule{\textwidth}{1.4pt}\par}
  \vspace{1.6cm}
  {\Huge\bfseries Text-to-SQL con QLoRA\par}
  \vspace{0.35cm}
  {\Large Fine-tuning de Qwen2.5-Coder para consultas MySQL en español\par}
  \vspace{1.2cm}
  \begin{tcolorbox}[colback=mist,colframe=line,boxrule=0.6pt,arc=2mm]
    \textbf{Objetivo.} Generar una única consulta MySQL 8.4.10 correcta a partir de un esquema y una pregunta en español, comparando el modelo base con dos configuraciones QLoRA mediante exactitud por ejecución.
  \end{tcolorbox}
  \vfill
  \begin{tabular}{@{}ll@{}}
    \textbf{Wilson Alejandro Calderón Argueta} & 22018 \\
    \textbf{Pablo Daniel Barillas Moreno} & 22193 \\
  \end{tabular}
  \vspace{0.8cm}

  Deep Learning y Sistemas Inteligentes · Sección 30\\
  2026
\end{titlepage}

\clearpage
\pagenumbering{roman}
\tableofcontents
\clearpage
\pagenumbering{arabic}

\section{Resumen}

% Sustituir únicamente cuando existan resultados reproducibles.
Este informe documenta el diseño del dataset, el baseline, dos configuraciones QLoRA que difieren solo en el rango, un experimento de inferencia guiado por ejecución y el análisis de errores. La exactitud final y la predicción de evaluación oculta se completarán a partir de los artefactos en \texttt{resultados/}.

\begin{finding}{Resultado principal}
  EX del baseline: \pending. EX del adapter seleccionado: \pending. Mejora absoluta: \pending.
\end{finding}

\section{Problema y restricciones}

Describir el contrato de \texttt{generar\_sql}, MySQL 8.4.10, límites de tiempo, determinismo y evaluación por ejecución. Explicar por qué la interfaz web queda fuera de la ruta evaluada.

\section{Datos}

\subsection{Fuentes y licencias}
Documentar Spider/BIRD, traducción, esquemas propios y condiciones de redistribución.

\subsection{Conversión y validación}
Reportar candidatos, aceptados y descartes por motivo. Toda cifra debe provenir de un manifiesto versionado.

\subsection{Separación por esquema}
Demostrar que no existe intersección entre entrenamiento y validación.

\section{Metodología}

\subsection{Baseline y prompt}
Incluir el prompt final y explicar que permanece idéntico con y sin adapter.

\subsection{Configuraciones QLoRA}

\begin{table}[H]
  \centering
  \caption{Comparación controlada de configuraciones.}
  \begin{tabularx}{\textwidth}{@{}lXXX@{}}
    \toprule
    Configuración & Rango & Factor constante & Hipótesis \\
    \midrule
    A & 16 & $\alpha=32$, dropout 0.05 & Menor costo y capacidad suficiente \\
    B & 32 & $\alpha=32$, dropout 0.05 & Mayor capacidad en consultas complejas \\
    \bottomrule
  \end{tabularx}
\end{table}

\subsection{Experimento de inferencia}
Comparar generación directa con una reparación ante error devuelto por MySQL.

\section{Resultados}

\subsection{Curvas y costo}
Insertar curvas train/eval loss, parámetros entrenables, tiempo y memoria.

\subsection{Exactitud por ejecución}
Reportar EX total y por nivel, tasa de ejecución y latencia p50/p95.

\section{Análisis de errores}
Clasificar sintaxis, entidad inexistente, JOIN, agregación, filtro y resultado distinto. Incluir ejemplos representativos sin seleccionar solo casos favorables.

\section{Predicción de evaluación oculta}
Construir la estimación desde la validación inédita, un intervalo bootstrap y una penalización por cambio de dominio.

\section{Conclusiones}
Responder qué aportó el fine-tuning, qué configuración ganó, cuándo ayudó la reparación y cuál es la principal limitación.

\section*{Uso de inteligencia artificial}
OpenAI Codex apoyó la planificación, estructura documental y handoff de interfaz; v0 de Vercel podrá apoyar la implementación React. El equipo revisará y validará todo artefacto. Ninguna API externa participa en la generación SQL evaluada.

\end{document}
